\section{Introducción: Contexto del Sector Astronómico}

\subsection{\textquestiondown Qué es el procesamiento de imágenes astronómicas?}

La astronomía observacional moderna depende de telescopios robóticos que capturan miles de imágenes cada noche. Estas imágenes contienen valiosa información científica, pero están afectadas por \textbf{ruido}: fluctuaciones aleatorias causadas por la atmósfera terrestre, detectores electrónicos y radiación cósmica.

El \textbf{denoising} (reducción de ruido) es el proceso de aplicar algoritmos computacionales para recuperar señales científicamente valiosas mientras se minimiza la introducción de artefactos (objetos falsos). Es una tarea crítica en que se toman decisiones fundamentales: \textquestiondown priorizar detectar todas las fuentes débiles (completitud) o priorizar la precisión de lo detectado (pureza)?

\subsection{El Observatorio Vera Rubin (LSST)}

El Observatorio Vera Rubin, con su Telescopio Sinóptico de Sondeo Amplio (LSST por sus siglas en inglés), está diseñado para revolucionar la astronomía mediante:

\begin{itemize}
    \item \textbf{Escala masiva}: Capturará 20 terabytes de datos por noche
    \item \textbf{Cobertura profunda}: Detectará galaxias débiles a distancias cosmológicas
    \item \textbf{Acceso público}: Los datos se publicarán sin restricción temporal
\end{itemize}

Sin embargo, la \textbf{participación científica} en proyectos LSST está limitada actualmente a instituciones miembros de la alianza oficial, concentradas en países desarrollados. Esta brecha excluye a investigadores de Latinoamérica, África y otras regiones del Sur Global.

\subsection{El Desafío Ético: Acceso Desigual}

El problema central que motiva este proyecto es la siguiente inequidad:

\textit{``Un investigador en Colombia o Perú que quiera validar o contribuir a métodos de procesamiento LSST enfrenta barreras técnicas, económicas e institucionales para acceder a pipelines oficiales y comunidades científicas de punta.''}

La visión del proyecto es establecer una \textbf{plataforma reproducible, extensible y centrada en métrica científica} que sea accesible sin restricción geográfica o institucional. Esto responde directamente a la inequidad descrita: mientras pipelines oficiales LSST están limitados a miembros de la alianza, Astro-Denoise propone una alternativa ética mediante código y datos públicos.

Este proyecto responde eticamente a esa exclusión:

\begin{quote}
\centering
\textbf{\textquestiondown Cómo podemos democratizar el acceso a herramientas científicas de punta manteniendo estándares rigurosos de reproducibilidad?}
\end{quote}

\subsection{Importancia Global}

La ciencia astronómica es patrimonio de la humanidad. Limitar el acceso a datos y métodos a instituciones privilegiadas:

\begin{enumerate}
    \item Reduce la diversidad de perspectivas científicas
    \item Perpetúa desigualdades educativas y económicas
    \item Desaprovecha talento de investigadores en regiones subrepresentadas
    \item Ralentiza avances científicos que benefician a toda la sociedad
\end{enumerate}

Este documento sostiene que \textbf{Astro-Denoise es una respuesta ética legítima} a estos desafíos.
